\documentclass[10pt,a4paper]{amsart}
\usepackage[margin=19mm]{geometry}
\usepackage{amsmath,amssymb,mathtools}
\usepackage[scr=rsfs,cal=euler]{mathalfa}
\usepackage{xfrac}
\usepackage[unicode,hidelinks,bookmarks=false]{hyperref}
\newcommand{\ie}{i.e.\ }
\newcommand{\cf}{cf.\ }
\newcommand{\resp}{respectively\ }
\newcommand{\Subsection}{\subsection}
\providecommand{\nobreakdash}{\nobreak\hbox{-}}
\setlength{\emergencystretch}{2em}
\multlinegap0pt
\usepackage{bm}
\usepackage{bbm}
\usepackage{dsfont}
\usepackage{xspace}
\usepackage{cancel}
\usepackage{mathtools}
\usepackage{cleveref}
\usepackage{amsthm}
\makeatletter
\@ifundefined{theorem}{\newtheorem{theorem}{Theorem}}{}
\@ifundefined{lemma}{\newtheorem{lemma}[theorem]{Lemma}}{}
\makeatother
\title[Deterministic Kalman filters for uncertain dynamical systems]{Deterministic Kalman filters for uncertain dynamical systems}
\author{Karl Kunisch, Jesper Schr\"oder}
\date{Source: arXiv:2506.00463v1 (CC BY 4.0)}
\begin{document}
\maketitle

\noindent\emph{Source and licence.} Deterministic Kalman filters for uncertain dynamical systems. Version arXiv:2506.00463v1 (CC BY 4.0), distributed under the Creative Commons Attribution 4.0 International licence, as recorded in the arXiv licence metadata for this exact version and shown on the abstract page https://arxiv.org/abs/2506.00463v1. Full rights notice: licenses/arxiv-2506.00463-CC-BY-4.0.txt.\par\smallskip

\noindent\textit{Editorial scope.} Counted unit: the Lemma whose source label is \texttt{lem: estJ\_sigma} in the article (source0.tex lines 1029--1043, source label \texttt{lem: estJ\_sigma}), delivered together with the article's own complete original proof of it (source0.tex lines 1045--1293, one \texttt{proof} environment of 249 lines). No mathematics is written, repaired or re-derived: statement and proof are the article's own text line for line. Required context retained uncounted: eq: OCP\_star\_cost (lines 901--912); eq: OCP\_star (lines 906--911); lem: est* (lines 975--989). Every definition, display and label of those lines is retained unchanged. Omitted from this package and not redistributed: 1--900, 912--974, 990--1028, 1044--1044, 1294--59980. Nothing inside a retained excerpt is omitted or altered: every retained body is delivered exactly as the article's lines are written, so no omission receipt of any kind accompanies this package. Citations: the retained excerpts carry no citation command at all, so no bibliography entry of the article is transcribed and no reference is redistributed. Reference adaptation: none was needed and none was made; every reference inside the delivered unit resolves inside it, so no unresolved reference or citation is produced. Printed slips kept literal: no printed word, symbol, hypothesis or number of the counted statement or of its proof was altered. Layout and class: the article's own class is \texttt{amsart} with packages including babel, inputenc, fontenc, lmodern, microtype, comment, amssymb, amsmath, amsthm, mathtools, mathrsfs, accents, etoolbox, siunitx; the wrapper is the lane's pinned \texttt{amsart} editorial wrapper at 10pt on A4, which restates the theorem-like environments the retained text needs and adds the article's own macro definitions that the retained text needs (none), with extra wrapper packages (bm, bbm, dsfont, xspace, cancel, mathtools, cleveref). The delivered numbering is the unit's own. Assets: no retained span contains a figure environment, a graphics-inclusion command, a PDF-page inclusion, a table or a TikZ picture; no image file is copied into this package, the package declares no asset dependency and no image file is redistributed with it. Licence: version arXiv:2506.00463v1 of this article is distributed under the Creative Commons Attribution 4.0 International licence, as recorded in the arXiv licence metadata for this exact version and shown on the abstract page https://arxiv.org/abs/2506.00463v1. A full rights notice is delivered at licenses/arxiv-2506.00463-CC-BY-4.0.txt. Characters: every retained excerpt is pure ASCII, so the delivered engine, XeLaTeX with its default Latin Modern text font, reports no missing character.
\begin{align}
    &\min_{x \in \mathcal{H}_t^T, u \in \mathcal{L}_t^T}~ \mathcal{J}_*(x,u;t,x_0) = 
    \frac{1}{2} \int_t^T \Vert u(s) \Vert_{Q_*}^2 
    + \Vert B^\top x(s) \Vert_{R_*}^2 \, \mathrm{d}s
    +\frac{1}{2} \Vert x(T) \Vert_{\Gamma_*}^2,\label{eq: OCP_star_cost} \\
    \begin{split}
    &\text{subject to}
    ~~~~~~~~~~\dot{x}(s) = A_*^\top x(s) + C^\top u(s)
    ~~~~~ s \in (t,T), \label{eq: OCP_star}\\
    &~~~~~~~~~~~~~~~~~~~~~~~x(t) = x_0,
    \end{split}
\end{align}

\begin{lemma}\label{lem: est*}
    Let $(\bar{x}_*,\bar{u}_*)$ be the unique solution to \eqref{eq: OCP_star_cost}-\eqref{eq: OCP_star} with associated adjoint state $p^*$. There exists $c_*>0$ independent of $t$ such that
    %
    \begin{equation*}
    \begin{aligned}
        \Vert \bar{x}_* \Vert_{\mathcal{H}_t^T}^2
        + \Vert \bar{x}_* \Vert_{\mathcal{C}_t^T}^2
        + \Vert \bar{u}_* \Vert_{\mathcal{L}_t^T}^2
        + \Vert p_* \Vert_{\mathcal{H}_t^T}^2
        + \Vert p_* \Vert_{\mathcal{C}_t^T}^2
        \leq c_* \Vert x_0 \Vert^2.
    \end{aligned}
    \end{equation*}
    %
\end{lemma}

\begin{lemma}\label{lem: estJ_sigma}
    Let $t \in [0,T)$ and $x_0 \in \mathbb{R}^n$ with $\Vert x_0 \Vert \leq 1$. Then there exists a constant $c_\mathcal{J} > 0 $ independent of $t$ such that 
    %
    \begin{equation*}
        \mathcal{J}_{\bar{\sigma}}
        ( \bar{x}_* - \bar{x}_{\bar{\sigma}},
        \bar{u}_* - \bar{u}_{\bar{\sigma}};
        t, x_0)
        \leq
        c_\mathcal{J}
        \Vert \mathcal{S}_* - \mathcal{S}_{\bar{\sigma}} \Vert_2^2
    \end{equation*}
    %
    where $(\bar{x}_*,\bar{u}_*)$ and $(\bar{x}_{\bar{\sigma}},\bar{u}_{\bar{\sigma}})$ are the optimal pairs associated with $\mathcal{J}_*$ and $\mathcal{J}_{\bar{\sigma}}$, respectively.
\end{lemma}

\begin{proof}
    We have $\nu_*(t,x_0) = \mathcal{J}_*(\bar{x}_*,\bar{u}_*;t,x_0)$, where $(\bar{x}_*,\bar{u}_*)$ is the unique minimizer of \eqref{eq: OCP_star_cost}-\eqref{eq: OCP_star} which is characterized via the optimality system
    %
    \begin{equation*}
    \begin{alignedat}{2}
        \dot{\bar{x}}_*(s) &= A_*^\top \bar{x}_*(s) + C^\top \bar{u}_*(s),~~~~~~~~~~
        &&\bar{x}_*(t) = x_0,\\
        %
        \dot{p}_*(s) &= - A_* p_*(s)- BR_*B^\top \bar{x}_*(s),
        &&p_*(T) = \Gamma_* \bar{x}_*(T),\\
        \bar{u}_*(s) &= - Q_*^{-1} C p_*(s), 
    \end{alignedat}
    \end{equation*}
    %
    where the time dependent equations hold for $s \in (t,T)$. 
    Analogously we have $\nu_{\bar{\sigma}}(t,x_0) = \mathcal{J}_{\bar{\sigma}}(\bar{x}_{\bar{\sigma}},\bar{u}_{\bar{\sigma}};t,x_0)$, where $(\bar{x}_{\bar{\sigma}},\bar{u}_{\bar{\sigma}})$ is the corresponding unique minimizer characterized via the optimality system
    %
    \begin{equation*}
    \begin{alignedat}{2}
        \dot{\bar{x}}_{\bar{\sigma}}(s) &= A_{\bar{\sigma}}^\top \bar{x}_{\bar{\sigma}}(s) + C^\top \bar{u}_{\bar{\sigma}}(s),~~~~~~~~~~
        &&\bar{x}_{\bar{\sigma}}(t) = x_0,\\
        %
        \dot{p}_{\bar{\sigma}}(s) &= - A_{\bar{\sigma}} p_{\bar{\sigma}}(s)- BR_{\bar{\sigma}}B^\top \bar{x}_{\bar{\sigma}}(s),
        &&p_{\bar{\sigma}}(T) = \Gamma_{\bar{\sigma}} \bar{x}_{\bar{\sigma}}(T),\\
        \bar{u}_{\bar{\sigma}}(s) &= - Q_{\bar{\sigma}}^{-1} C p_{\bar{\sigma}}(s).
    \end{alignedat}
    \end{equation*}
    %
    We define
    %
    \begin{equation*}
    \begin{alignedat}{3}
        x_{\delta} &\coloneqq \bar{x}_* - \bar{x}_{\bar{\sigma}},
        ~~~~~
        u_{\delta} &&\coloneqq \bar{u}_* - \bar{u}_{\bar{\sigma}},
        ~~~~~~~~~~
        p_{\delta} &&\coloneqq \bar{p}_* - \bar{p}_{\bar{\sigma}},\\
        %
        \Gamma_{\delta} & \coloneqq \Gamma_* - \Gamma_{\bar{\sigma}},
        ~~~~~
        R_{\delta} &&\coloneqq R_* - R_{\bar{\sigma}},
        ~~~~~
        Q_{\delta^{-1}} &&\coloneqq Q_*^{-1} - Q_{\bar{\sigma}}^{-1},
        ~~~~~
        A_{\delta} = A_* - A_{\bar{\sigma}},
    \end{alignedat}
    \end{equation*}
    %
    and obtain
    %
    \begin{equation*}
    \begin{alignedat}{2}
        \dot{x}_{\delta} 
        &= A_{\bar{\sigma}}^\top x_{\delta}
            + A_{\delta}^\top \bar{x}_*
            + C^\top u_{\delta},
        ~~~~~
        ~~~~~
        ~~~~~
        ~~~~~
        ~~~~~
        ~~~~
        x_{\delta} (t) &&= 0,\\
        \dot{p}_{\delta}
        &= - A_{\bar{\sigma}} p_{\delta} 
            - A_{\delta} p_* 
            - B R_{\bar{\sigma}} B^\top x_{\delta}
            - B R_{\delta} B^\top \bar{x}_*,
        ~~~~~
        p_{\delta} (T) 
        &&= \Gamma_{\bar{\sigma}} x_{\delta} (T)
            + \Gamma_{\delta} \bar{x}_*(T),\\
        u_{\delta} 
        &= - Q_{\bar{\sigma}}^{-1} C p_{\delta} 
            - Q_{\delta^{-1}} C p_*.
    \end{alignedat}
    \end{equation*}
    %
    Now a combination of standard calculus, Gronwall's Lemma, the norm bound on $x_0$, and \Cref{lem: est*} yield existence of a constant $c_\delta>0$ depending on $T$, $c_*$, $\Vert A_{\bar{\sigma}} \Vert$, $\Vert B \Vert$, $\Vert C\Vert$, $\Vert R_{\bar{\sigma}} \Vert$, and $\Vert \Gamma_{\bar{\sigma}} \Vert$ but independent of $t$ such that
    %
    \begin{equation}\label{eq: est_x_p_delta}
    \begin{aligned}
        \Vert x_{\delta} \Vert_{\mathcal{C}_t^T}^2
        +
    	\Vert x_{\delta} \Vert_{\mathcal{H}_t^T}^2
    	&\leq c_\delta \left( \Vert A_\delta \Vert^2 
    	+ 
    	\Vert u_\delta \Vert_{\mathcal{L}_t^T}^2 \right),\\
    	%
        \Vert p_\delta \Vert_{\mathcal{C}_t^T}^2
        +
    	\Vert p_\delta \Vert_{\mathcal{H}_t^T}^2	
    	&\leq c_\delta \left(
    	\Vert u_\delta \Vert_{\mathcal{L}_t^T}^2
    	+ 
    	\Vert A_\delta \Vert^2 
    	+ 
    	\Vert R_\delta \Vert^2 
    	+
    	\Vert \Gamma_\delta \Vert^2 
    	\right).
    \end{aligned}
    \end{equation}
    %
    
    Further, by testing the equation for $\dot{p}_{\delta}$ with $x_{\delta}$ we find 
    %
    \begin{equation}\label{eq: dual1}
        -\int_t^T (\dot{p}_{\delta}, x_{\delta}) \, \mathrm{d}s
        =
        \int_t^T
        (A_{\bar{\sigma}} p_{\delta}, x_{\delta} )
        + 
        (A_{\delta} p_*, x_{\delta} )
        +
        (B R_{\bar{\sigma}} B^\top x_{\delta}, x_{\delta} )
        +
        (B R_{\delta} B^\top \bar{x}_*, x_{\delta})
        \, \mathrm{d}s.
    \end{equation}
    %
    On the other hand, with $x_{\delta} (0) = 0$ and the identity for $p_{\delta} (T)$ we find testing the equation for $\dot x_\delta$ with $p_\delta$
     %
    \begin{equation}\label{eq: dual2}
    \begin{aligned}
        &-\int_t^T (\dot{p}_{\delta}, x_{\delta}) \, \mathrm{d}s
        = \int_t^T (\dot{x}_{\delta}, p_{\delta}) \, \mathrm{d}s
            - (p_{\delta}(T),x_{\delta}(T))\\
        &= \int_t^T 
        (A_{\bar{\sigma}}^\top x_{\delta}, p_{\delta})
        + (A_{\delta}^\top \bar{x}_*, p_{\delta} )
        + ( C^\top u_{\delta}, p_{\delta} )
        \, \mathrm{d}s
        - ( \Gamma_{\bar{\sigma}} x_{\delta}(T), x_{\delta}(T) )
        - ( \Gamma_{\delta} \bar{x}_*, x_{\delta}(T) ).
    \end{aligned}
    \end{equation}
    %
    By subtracting \eqref{eq: dual1} from \eqref{eq: dual2}, rearranging terms and utilizing the identity for $u_{\delta}$ we arrive at
    %
    \begin{equation*}
    \begin{aligned}
        \delta \mathcal{J}_{\bar{\sigma}}
        &\coloneqq
        \int_t^T 
        \Vert u_{\delta} \Vert_{Q_{\bar{\sigma}}}^2
        + \Vert B^\top x_{\delta} \Vert_{R_{\bar{\sigma}}}^2
        \, \mathrm{d}s
        + \Vert x_{\delta}(T) \Vert_{\Gamma_{\bar{\sigma}}}^2\\
        &=
        - \int_t^T  
        (A_{\delta} p_* , x_{\delta})
        + (B R_{\delta} B^\top \bar{x}_*, x_{\delta})
        - ( A_{\delta}^\top \bar{x}_*, p_{\delta} )
        + ( Q_{\bar{\sigma}} u_{\delta}, Q_{\delta^{-1}} p_* )
        \, \mathrm{d}s\\
        &- ( \Gamma_{\delta} \bar{x}_*(T), x_{\delta}(T) ).
    \end{aligned}
    \end{equation*}
    %
    Utilizing \Cref{lem: est*}, \eqref{eq: est_x_p_delta}, and Young's inequality it follows that for any $\epsilon>0$ it holds
    %
    \begin{equation*}
    \begin{aligned}
    	\delta \mathcal{J}_{\bar{\sigma}}
    	&\leq
    	c_* \int_t^T
    		\Vert A_\delta \Vert \Vert x_\delta \Vert
    		+
    		\Vert B \Vert^2 \Vert R_\delta \Vert \Vert x_\delta \Vert
    		+
    		\Vert A_\delta \Vert  \Vert p_\delta \Vert
    		+
    		  \Vert Q_{\bar{\sigma}} \Vert \Vert Q_{\delta^{-1}} \Vert \Vert u_\delta \Vert 
     	\, \mathrm{d}s\\
     	&+
        c_* \Vert \Gamma_\delta \Vert \Vert x_\delta \Vert_{\mathcal{C}_t^T}\\
     	%
     	&\leq
     	c_* \int_t^T
     	\frac{1}{\epsilon} \Vert A_\delta \Vert^2
     	+
     	\epsilon \Vert x_\delta \Vert^2
     	+
     	\frac{\Vert B \Vert^4}{2\epsilon} \Vert R_\delta \Vert^2
     	+
     	\frac{\epsilon}{2} \Vert p_\delta \Vert^2 
     	+
     	\frac{\Vert Q_{\bar{\sigma}} \Vert^2}{2 \epsilon} \Vert Q_{\delta^{-1}} \Vert^2
     	+ 
     	\frac{c_* \epsilon}{2} \Vert u_\delta \Vert^2
     	\, \mathrm{d}s\\
     	&+ 
     	\frac{c_*}{2\epsilon} \Vert \Gamma_\delta \Vert^2
     	+   
     	\frac{\epsilon}{2} \Vert x_\delta \Vert_{\mathcal{C}_t^T}^2\\
     	&\leq
     	\frac{c_* T}{\epsilon} \Vert A_\delta \Vert^2
     	+
     	\frac{c_* T \Vert B \Vert^4}{2\epsilon} \Vert R_\delta \Vert^2
     	+
     	\frac{c_* T \Vert Q_{\bar{\sigma}} \Vert^2}{2\epsilon} \Vert Q_{\delta^{-1}} \Vert^2
     	+
     	\frac{c_*}{2\epsilon} \Vert \Gamma_\delta \Vert^2
        + 
     	\frac{c_* \epsilon}{2} \Vert u_\delta \Vert_{\mathcal{L}_t^T}^2\\
     	&+ 
     	(\epsilon c_* c_\delta + \frac{\epsilon c_\delta}{2}) (\Vert A_\delta \Vert^2 + \Vert u_\delta \Vert_{\mathcal{L}_t^T}^2)
     	+ \frac{\epsilon c_* c_\delta}{2} \left(
     	\Vert u_\delta \Vert_{\mathcal{L}_t^T}^2
     	+ 
     	\Vert A_\delta \Vert^2 
     	+ 
     	\Vert R_\delta \Vert^2 
     	+
     	\Vert \Gamma_\delta \Vert^2 
     	\right).
	\end{aligned}  
    \end{equation*}
    %
    Since 
    %
    \begin{equation*}
        \Vert u_\delta \Vert_{\mathcal{L}_t^T}^2 \leq \Vert Q_{\bar{\sigma}}^{-\tfrac{1}{2}} \Vert^2 \int_t^T \Vert u_\delta \Vert_{Q_{\bar{\sigma}}}^2 \, \mathrm{d}s,
    \end{equation*}
    %
    for a sufficiently small $\epsilon > 0$ all terms on the right hand side depending on $u_\delta$ can be absorbed in the left hand side, and with appropriate constants $c_1, c_2 > 0$ it follows that 
    %
    \begin{equation*}
        c_1 \delta \mathcal{J}_{\bar{\sigma}}
        \leq 
        c_2 \left( 
        \Vert A_\delta \Vert + \Vert \Gamma_\delta \Vert + \Vert R_\delta \Vert + \Vert Q_{\delta^{-1}} \Vert
        \right).
    \end{equation*}
    %
    Finally, note that the mapping $ \mathcal{F} (A) = A^{-1}$ is continuously differentiable on
    the open set of positive definite matrices with derivative $D\mathcal{F}(A)B = - A^{-1}B A^{-1}$. Hence an application of Taylor's theorem yields
    %
    \begin{equation*}
        \Vert Q_{\delta^{-1}} \Vert
        =
        \Vert Q_*^{-1} - Q_{\bar{\sigma}}^{-1} \Vert
        \leq 
        \Vert Q_{\bar{\sigma}}^{-1} \Vert^2 \Vert Q_* - Q_{\bar{\sigma}} \Vert
    \end{equation*}
    %
    and the assertion follows.
\end{proof}

\end{document}
